\textbf{Definitions.}
\begin{itemize}
\item $\tau(n)$ denotes the number of positive divisors of $n\in\mathbb{N}$.
\item Let $f:\mathbb{N}\to\mathbb{R}$ and define $k_n=\lfloor f(n)\rfloor\in\mathbb{Z}_{\ge 0}$.
\item For $n\in\mathbb{N}$, define
\[
F(f,n)=\frac{\tau\big((n+k_n)!\big)}{\tau(n!)}.
\]
\item Throughout, $\log$ denotes the natural logarithm.
\end{itemize}

\textbf{Questions.}
\begin{itemize}
\item[(Q1)] For sufficiently large constants $C>0$, is it true that
\[
\lim_{n\to\infty} F\big((\log n)^C,n\big)=\infty?
\]
\item[(Q2)] Is the set $\{F(\log n,n): n\in\mathbb{N}, n\ge 3\}$ dense in $(1,\infty)$?
\item[(Q3)] More generally, if $f(n)\le \log n$ is monotone and $f(n)\to\infty$, is the set $\{F(f,n):n\in\mathbb{N}\}$ dense in $(1,\infty)$?
\end{itemize}
